\textbf{\textit{Soal. }}In triangle $ABC,AB>AC,I$ is the incenter, $AM$ is the midline. The line crosses $I$ and is perpendicular to $BC $ intersect $AM$ at point $L$, and the symmetry of $I$ with respect to point $A$ is $J$
Prove that: $\angle ABJ= \angle LBI$.

\begin{proof}[\textbf{Solusi. }] (Sunday, 30 October 2022) Misalkan $N$ adalah proyeksi $I$ pada garis $AC$, $I_A$ adalah titik pusat $A-$excircle $\triangle ABC$, $K$ adalah proyeksi $I_A$ pada $BC$, dan $T$ adalah titik pada garis $II_A$ sedemikian sehingga $TM \perp BC$. 

Perhatikan karena $\angle NIA = \angle CIA - \angle CIN = (90^\circ+\frac{\angle B}{2})-90^\circ=\frac{\angle B}{2} = \angle IBA$ dan $\angle NAI = \angle IAB$ maka didapat $\triangle BIN \sim \triangle BAJ$.

Selanjutnya, observasi bahwa
$$\angle NIB = 360^\circ - \angle BIC - \angle CIN
            = 360^\circ - \left(90^\circ + \dfrac{\angle A}{2}\right) - 90^\circ
            = 180^\circ - \dfrac{1}{2}\angle IAB
            = \angle JAB.$$
Karena $\triangle BIN \sim \triangle BAJ$ maka $\dfrac{NI}{AJ}=\dfrac{NI}{IA}=\dfrac{IB}{AB}$. Hal ini menyebabkan $\triangle BIN \sim \triangle BAJ$ yang secara langsung menunjukkan bahwa $\angle JBA = \angle NBI$. Oleh karena itu, untuk membuktikan $\angle JBA = \angle LBI$ cukup dibuktikan $N,L,B$ segaris.
% \begin{center}
    \begin{asy}
        import olympiad;
        import geometry;
        size(350);

        pair A, B, C;
        A = dir(60);
        B = dir(-27);
        C = dir(-153);
        draw(A--B--C--A);
        dot("$A$", A, NE);
        dot("$B$", B, SE);
        dot("$C$", C, SW);

        pair I = incenter(triangle(A,B,C));
        dot("$I$", I, NE);

        pair M = (B+C)/2;
        dot("$M$", M, SW);
        draw(A--M);

        pair D = foot(I, B, C);
        pair L = intersectionpoint(line(A,M),line(D,I));
        pair J = rotate(180, A)*I;
        dot("$D$", D, SE);
        dot("$L$", L, NW);
        dot("$J$", J, NW);
        draw(D--L);
        draw(B--I);
        draw(C--I);

        pair T = intersectionpoint(line(A,I), perpendicular(M, line(B,C)));
        dot("$T$", T, SE);
        draw(M--T);

        pair Ia = rotate(180, T)*I;
        dot("$I_A$", Ia, NW);
        draw(J--A--I--Ia);

        pair N = intersectionpoint(line(A,C),perpendicular(I,line(C,I)));
        dot("$N$", N, NW);
        draw(I--N, blue);
        draw(rightanglemark(C,I,N,2), blue);
        draw(rightanglemark(T,M,B,2));
        draw(rightanglemark(I,D,B,2));

        pair K = foot(Ia, C, B);
        dot("$K$", K, SW);
        draw(Ia--K);
        draw(C--K, red);
        draw(B--D, red);
        draw(rightanglemark(Ia,K,B,2));
        
        draw(B--J--A--B, blue);
        draw(B--I--N--B, blue);
        draw(C--Ia, blue);
    \end{asy}
\end{center}
\begin{center}
    \begin{tikzpicture}[scale=3, dot/.style={circle, fill, inner sep=0.5pt, outer sep=0pt}]

        \tkzDefPoint(60:1){A}
        \tkzDefPoint(-27:1){B}
        \tkzDefPoint(-153:1){C}

        \tkzInCenter(A,B,C)\tkzGetPoint{I}
        \tkzDefMidPoint(B,C)\tkzGetPoint{M}
        \tkzDefPointBy[projection=onto B--C](I)\tkzGetPoint{D}
        \tkzInterLL(A,M)(D,I)\tkzGetPoint{L}
        \tkzDefPointBy[symmetry=center A](I)\tkzGetPoint{J}

        % T perpotongan AI dengan tegak lurus BC melalui M
        \tkzDefLine[orthogonal=through M](B,C)\tkzGetPoint{mp}
        \tkzInterLL(A,I)(M,mp)\tkzGetPoint{T}
        \tkzDefPointBy[symmetry=center T](I)\tkzGetPoint{Ia}

        % N perpotongan AC dengan tegak lurus CI melalui I
        \tkzDefLine[orthogonal=through I](C,I)\tkzGetPoint{cip}
        \tkzInterLL(A,C)(I,cip)\tkzGetPoint{N}

        \tkzDefPointBy[projection=onto C--B](Ia)\tkzGetPoint{K}

        \tkzDrawPolygon(A,B,C)
        \tkzDrawSegment(A,M)
        \tkzDrawSegment(D,L)
        \tkzDrawSegments(B,I C,I)
        \tkzDrawSegment(M,T)
        \tkzDrawSegments(J,A A,I I,Ia)
        \tkzDrawSegment(Ia,K)
        \tkzDrawSegment[red](C,K)
        \tkzDrawSegment[red](B,D)
        \tkzDrawSegments[blue](I,N B,J J,A A,B B,I N,B C,Ia)

        \tkzMarkRightAngle[blue,size=0.16](C,I,N)
        \tkzMarkRightAngle[size=0.16](T,M,B)
        \tkzMarkRightAngle[size=0.16](I,D,B)
        \tkzMarkRightAngle[size=0.16](Ia,K,B)

        \tkzDrawPoints(A,B,C,I,M,D,L,J,T,Ia,N,K)
        \tkzLabelPoint[above right](A){$A$}
        \tkzLabelPoint[above right](B){$B$}
        \tkzLabelPoint[above left](C){$C$}
        \tkzLabelPoint[above right](I){$I$}
        \tkzLabelPoint[below left](M){$M$}
        \tkzLabelPoint[below right](D){$D$}
        \tkzLabelPoint[above left](L){$L$}
        \tkzLabelPoint[above left](J){$J$}
        \tkzLabelPoint[below right](T){$T$}
        \tkzLabelPoint[above left](Ia){$I_A$}
        \tkzLabelPoint[above left](N){$N$}
        \tkzLabelPoint[below left](K){$K$}

    \end{tikzpicture}
\end{center}

Sekarang perhatikan bahwa $\angle INA = 180^\circ - \angle CNI = 90^\circ + \angle ICN = \angle I_ACI + \angle ICN = \angle I_ACA$ yang mengakibatkan $NI \parallel CI_A$. Diperoleh $\dfrac{II_A}{IA}=\dfrac{CN}{AN}$.

Di lain pihak, dengan properti terkenal, kita punya $CK=DB$ dan $KM=MD$. Hal ini menyebabkan $II_A = 2 IT$. Lalu, karena $MT, LI \perp BC$ maka $MT \parallel LI$ sehingga $\dfrac{IT}{IA}=\dfrac{ML}{AL}$. Dari sini kita peroleh
\begin{align*}
    \dfrac{CN}{AN} = \dfrac{II_A}{IA} = \dfrac{2IT}{IA} = \dfrac{2ML}{AL}
\end{align*}
yang mengakibatkan
\begin{align*}
    \dfrac{AB}{AN} \cdot \dfrac{CN}{CB} = \dfrac{AB}{AL} \cdot \dfrac{ML}{\frac{1}{2}CB} = \dfrac{AB}{AL} \cdot \dfrac{ML}{MB}.
\end{align*}
Padahal dengan dalil sinus pada $\triangle ABN$ dan $\triangle NBC$ serta $\triangle ABL$ dan $\triangle LBM$, akan didapat
\begin{align*}
    \dfrac{AB}{AN}\cdot\dfrac{CN}{CB} = \dfrac{\sin \angle BNA}{\sin \angle ABN} \cdot \dfrac{\sin \angle NBC}{\sin \angle CNB} = \dfrac{\sin \angle NBC}{\sin \angle ABN}
\end{align*}
dan
\begin{align*}
    \dfrac{AB}{AL}\cdot\dfrac{ML}{MB} = \dfrac{\sin \angle BLA}{\sin \angle ABL} \cdot \dfrac{\sin \angle LBM}{\sin \angle MLB} = \dfrac{\sin \angle LBM}{\sin \angle ABL}
\end{align*}
sehingga sekarang kita punya
\begin{align*}
    \dfrac{\sin \angle NBC}{\sin \angle ABN} = \dfrac{\sin \angle LBM}{\sin \angle ABL}.
\end{align*}
Namun, kita mengetahui bahwa $\angle ABN + \angle NBC = \angle ABL + \angle LBC = \angle ABC < 180^\circ$. Hal ini memaksa $\angle ABN = \angle ABL$ dan $\angle NBC = \angle LBC$ atau setara dengan titik-titik $N,L,B$ segaris, seperti yang ingin dibuktikan.
\end{proof}